\section{Background}
\label{sec:background}

\subsection{Threat model}
The attacker controls a share of the training data and plants a trigger with a target label. The defender receives the trained model and holds a small clean validation set drawn from the same distribution, here \HeldoutSize{} images, but never the training set and never a poisoned example. At deployment the defender sees 1 input at a time and must decide whether it carries a trigger before returning a prediction. This is the setting of the input-level detectors of \cref{sec:related}.

\subsection{The prediction shift phenomenon in backdoored networks}
PSBD \cite{li2025psbd} runs the model $k$ times on an input $x$ with dropout at rate $p$ active and compares each dropout prediction with the prediction without dropout. The paper defines the prediction shift indicator and the shift ratio of a set $\mathcal{D}$ as
\begin{equation}
  \begin{split}
    \ps(x) &= \mathbb{I}\big[\mathcal{Y}(x;\theta) \neq \mathcal{Y}(x;\theta')\big], \\
    \sigma(\mathcal{D}) &= \frac{1}{k|\mathcal{D}|}\sum_{x\in\mathcal{D}} \ps(x),
  \end{split}
\end{equation}
where $\theta$ are the weights, $\theta'$ the weights with dropout active and $\mathcal{Y}$ the predicted class. The detection statistic, the prediction shift uncertainty, is the drop in the probability of the no-dropout class averaged over the $k$ passes,
\begin{equation}
  \psu(x) = \frac{1}{k}\sum_{i=1}^{k}
    \big[ P(\hat y \mid x;\theta) - P(\hat y \mid x;\theta'_i) \big],
\end{equation}
with $\hat y$ the class predicted without dropout. A poisoned input keeps its probability under dropout and has a low $\psu$.

The dropout rate is chosen so that the shift ratio of the clean validation set reaches \AdaptiveShiftTarget{}. The threshold is the \HeadlineQuantilePercent{} percentile of the validation set's $\psu$ at $k = 3$ passes. In a ResNet the paper applies dropout after each residual connection of the basic block, and it explains the phenomenon by neuron bias, which \cref{sec:mechanism:phenomenon} tests on ViT.

\subsection{Perturbation sites in a ViT block}
A ViT-B/16 \cite{dosovitskiy2021an} splits a \ModelInputSize{} by \ModelInputSize{} image into \VitPatches{} patches of 16 by 16 pixels, embeds each as a \VitWidth{}-dimensional token, prepends a class token and runs \VitBlocks{} identical blocks. Each block is pre-norm \cite{xiong2020layernorm}: it applies a LayerNorm, multi-head self-attention over the \VitTokens{} tokens and a residual add, then a second LayerNorm, a 2-layer MLP applied to every token independently and a second residual add. The classifier reads the class token after a final LayerNorm. \Cref{fig:vit-block} shows 1 block with the sites where we inject a perturbation.

Attention is the only operation that mixes information across tokens. A perturbation before the attention LayerNorm therefore acts on each token before any mixing, a perturbation before the MLP LayerNorm acts after the block's mixing has happened, and a perturbation on the residual stream acts on the sum of everything the network has written so far.

\begin{figure}[t]
\centering
\begin{tikzpicture}[
  node distance=3mm,
  box/.style={draw, rounded corners=1pt, minimum width=22mm, minimum height=6mm, font=\scriptsize, align=center},
  norm/.style={box, fill=gray!12},
  op/.style={box, fill=blue!8},
  plus/.style={draw, circle, inner sep=1pt, font=\scriptsize},
  site/.style={draw, circle, fill=orange!80, inner sep=1.4pt, font=\bfseries\tiny},
  flow/.style={-{Latex[length=1.6mm]}, thick},
  lbl/.style={font=\scriptsize, text width=34mm, align=left}
]
\node[font=\scriptsize] (in) {residual stream $x$ \quad (197 tokens $\times$ 768)};
\node[norm, below=of in] (ln1) {LayerNorm};
\node[op, below=of ln1] (attn) {multi-head attention};
\node[plus, below=of attn] (add1) {$+$};
\node[norm, below=of add1] (ln2) {LayerNorm};
\node[op, below=of ln2] (mlp) {MLP (per token)};
\node[plus, below=of mlp] (add2) {$+$};
\node[font=\scriptsize, below=of add2] (out) {to the next block};
\draw[flow] (in) -- (ln1); \draw[flow] (ln1) -- (attn); \draw[flow] (attn) -- (add1);
\draw[flow] (add1) -- (ln2); \draw[flow] (ln2) -- (mlp); \draw[flow] (mlp) -- (add2); \draw[flow] (add2) -- (out);
\draw[flow] (in.south) ++(0,-1mm) -| ++(-16mm,0) |- (add1.west);
\draw[flow] (add1.south) ++(0,-1mm) -| ++(-16mm,0) |- (add2.west);
\node[site, right=2mm of ln1.north east, yshift=1mm] (A) {A};
\node[site, right=2mm of ln2.north east, yshift=1mm] (B) {B};
\node[site, right=2mm of attn.south east, yshift=-1mm] (C) {C};
\node[site, right=2mm of add1.east, xshift=4mm] (D) {D};
\node[lbl, right=10mm of A] {A \textbf{before the attention norm}: token mask (recommended), channel mask, Gaussian noise};
\node[lbl, right=10mm of C] {C \textbf{before the residual add}: dropout on the branch output};
\node[lbl, right=10mm of D] {D \textbf{after the residual add}: dropout on the stream (the published placement)};
\node[lbl, right=10mm of B] {B \textbf{before the MLP norm}: token mask, Gaussian noise};
\end{tikzpicture}

\caption{1 ViT block with the sites we perturb, in forward order. Site A is the attention input, where PSBD-TM masks whole tokens. Site C is the residual stream after the add, where PSBD-RD applies dropout. The attention branch output before the add (B), the MLP input (D), the attention heads and the MLP hidden units are the other sites in the basis.}
\label{fig:vit-block}
\end{figure}
